% ============================================================
% f2014_tianjin_arts_q10_views.tex
% 2014 天津卷（文科）第 10 题：三视图
% 来源 PDF：普通高考/2014/2014天津文.pdf
% 原页：第 1 页（A3），印面 925，第 10 题题干下方
% 正视图、侧视图相同：宽 2、高 4 的矩形上加底宽 4、高 2 的三角形；
% 俯视图为半径 2 的圆，内含半径 1 的虚线圆（下部圆柱被遮挡的轮廓）。
% ============================================================
\begin{tikzpicture}[>=Stealth, line width=0.6pt, font=\small, scale=0.55,
  line cap=round, line join=round]
  % ---------- 正视图 ----------
  % 下部矩形：x∈[1,3]，高 4；上部三角形：底 x∈[0,4] 于 y=4，顶点 (2,6)
  \draw (1,0) -- (3,0) -- (3,4) -- (4,4) -- (2,6) -- (0,4) -- (1,4) -- (1,0);
  % 底宽 2（矩形）
  \draw[<->] (1,-0.5) -- (3,-0.5);
  \node[below] at (2,-0.5) {$2$};
  % 底宽 4（三角形底边）
  \draw[<->] (-0.6,-1.3) -- (4.6,-1.3);
  \draw (-0.6,-0.9) -- (-0.6,-1.7);
  \draw (4.6,-0.9) -- (4.6,-1.7);
  \node[below] at (2,-1.3) {$4$};
  \node[below] at (2,-2.2) {正视图};
  % 柱高 4、锥高 2（置于右侧）
  \draw[<->] (5.4,0) -- (5.4,4);
  \draw (5.0,0) -- (5.8,0);
  \draw (5.0,4) -- (5.8,4);
  \node[right] at (5.4,2) {$4$};
  \draw[<->] (5.4,4) -- (5.4,6);
  \draw (5.0,6) -- (5.8,6);
  \node[right] at (5.4,5) {$2$};

  % ---------- 侧视图 ----------
  \begin{scope}[shift={(9.2,0)}]
    \draw (1,0) -- (3,0) -- (3,4) -- (4,4) -- (2,6) -- (0,4) -- (1,4) -- (1,0);
    \draw[<->] (1,-0.5) -- (3,-0.5);
    \node[below] at (2,-0.5) {$2$};
    \draw[<->] (-0.6,-1.3) -- (4.6,-1.3);
    \draw (-0.6,-0.9) -- (-0.6,-1.7);
    \draw (4.6,-0.9) -- (4.6,-1.7);
    \node[below] at (2,-1.3) {$4$};
    \node[below] at (2,-2.2) {侧视图};
  \end{scope}

  % ---------- 俯视图 ----------
  \begin{scope}[shift={(2,-6.6)}]
    \draw (0,0) circle (2);
    \draw[dashed] (0,0) circle (1);
    \node[below] at (0,-2.7) {俯视图};
  \end{scope}
\end{tikzpicture}
